% Part: computability
% Chapter: tm-computations
% Section: church-turing-thesis

\documentclass[../../../include/open-logic-section]{subfiles}

\begin{document}

\olfileid[pt]{tur}{mac}{ctt}
\olsection{A Tese de Church--Turing}

As máquinas de Turing devem ser um substituto preciso do
conceito de procedimento efetivo. Turing pensava que quem
compreendesse tanto o conceito de procedimento efetivo quanto o de
máquina de Turing teria a intuição de que tudo o que pudesse
ser feito por meio de um procedimento efetivo poderia ser feito por uma máquina de Turing.
Essa afirmação é sustentada pelo fato de que todas as outras propostas
precisas de substituto do conceito de procedimento efetivo se mostram
extensionalmente equivalentes ao conceito de máquina de Turing
---ou seja, podem computar exatamente o mesmo conjunto de funções. Essa
afirmação chama-se \emph{tese de Church--Turing}.

\begin{defn}[Tese de Church--Turing]
A \emph{tese de Church--Turing} afirma que tudo o que é computável por meio
de um procedimento efetivo é computável por máquina de Turing.
\end{defn}

A tese de Church--Turing é invocada de duas maneiras. O primeiro tipo
de uso é uma desculpa para a preguiça. Suponha
que temos uma descrição de um procedimento efetivo para computar algo,
digamos, em ``pseudocódigo.'' Então podemos invocar a tese de Church--Turing
para justificar a afirmação de que a mesma função é computada por alguma máquina de Turing,
mesmo que não a tenhamos de fato construído.

O outro uso da tese de Church--Turing é mais filosoficamente
interessante. Pode-se mostrar que existem funções que não podem ser
computadas por máquinas de Turing. Disso, usando a tese de Church--Turing,
conclui-se que não podem ser computadas efetivamente, usando
qualquer procedimento que seja. Pois se houvesse tal procedimento, pela
tese de Church--Turing seguiria que haveria uma máquina de Turing para ele. Assim, se podemos provar que não há máquina de Turing
que o compute, também não pode haver um procedimento efetivo. Em
particular, a tese de Church--Turing é invocada para afirmar que o
chamado problema da parada não só não pode ser resolvido por máquinas de Turing,
como não pode ser resolvido efetivamente de modo algum.

\end{document}
